\section{账本事件的制度与社会深描：an-lushan-late-tang}

\label{sec:ledger-depth-an-lushan-late-tang}

本组以账本所列事件为可核锚点，转而追问军事、财政、地方行政、宗教与跨区域网络怎样在具体空间中相互限制。官方叙述、文书和物质遗存各有覆盖范围，不能由任何一种来源替代普通家庭与流动人群的完整经验。

\subsection{区域制度与资源}

条目\texttt{LT-755-119}所记的天宝十四载（755）十二月“安禄山军陷东都洛阳，叛军以洛阳为前进基地”，发生在唐与安禄山叛军的政治语境中。账本将其原因概括为范阳军南下速度超过唐廷临时调兵，河北、河东防线未能在洛阳以北形成有效阻截，并提示叛军获得连接关东与河南的枢纽，唐廷被迫以潼关和长安作为下一道防线。这类节点进入地方社会时，要经过官员、军队、书吏、商人、寺院、道路与家庭等不同中介；同一政策或战事在城市、边地和乡村不会具有相同的节奏。

核心来源为《旧唐书·玄宗本纪下》；《旧唐书·安禄山传》；《资治通鉴》唐纪（天宝十四载十二月条）。它们可以较有把握地支持年序、制度公布、使节活动或特定地点的记载，却不能自动给出人口、财富、兵额、族属和内心动机的完整调查。账本的证据注记是“洛阳失陷的月日与入城经过在本纪、列传间有异，唐方材料对守军和城内反应保存有限”。因此，不能将官府的分类、战报的数字、宗教文本的愿望或后出的总括直接写成全体社会事实。

更合适的解释是，将这一事件看成资源、信息和权力重新连接的一环：命令是否抵达，取决于交通和地方协作；征发是否落实，取决于登记、市场和劳力；文化和宗教网络能否延续，也受战争、迁徙和赞助条件影响。现有材料只照亮其中一部分，叙事必须让区域差异与材料沉默继续可见。

条目\texttt{LT-756-001}所记的天宝十五载（756）六月“哥舒翰出潼关战败，燕军西进，唐玄宗离长安入蜀”，发生在唐与燕的政治语境中。账本将其原因概括为洛阳失守后关中朝廷急于收复东都，杨国忠与哥舒翰的军政冲突使潼关守军被迫出战，并提示关中门户打开，皇帝出奔和京师失守使唐朝合法性与行政中枢同时遭受危机。这类节点进入地方社会时，要经过官员、军队、书吏、商人、寺院、道路与家庭等不同中介；同一政策或战事在城市、边地和乡村不会具有相同的节奏。

核心来源为《旧唐书·肃宗本纪》；《旧唐书·哥舒翰传》；《资治通鉴》唐纪（天宝十五载六月条）。它们可以较有把握地支持年序、制度公布、使节活动或特定地点的记载，却不能自动给出人口、财富、兵额、族属和内心动机的完整调查。账本的证据注记是“潼关败因常被归为宰相逼战，哥舒翰、杨国忠和玄宗的决策次序仍须比对不同叙事”。因此，不能将官府的分类、战报的数字、宗教文本的愿望或后出的总括直接写成全体社会事实。

更合适的解释是，将这一事件看成资源、信息和权力重新连接的一环：命令是否抵达，取决于交通和地方协作；征发是否落实，取决于登记、市场和劳力；文化和宗教网络能否延续，也受战争、迁徙和赞助条件影响。现有材料只照亮其中一部分，叙事必须让区域差异与材料沉默继续可见。

条目\texttt{LT-756-002}所记的至德元载（756）七月“太子李亨在灵武即位，遥尊玄宗为太上皇”，发生在唐的政治语境中。账本将其原因概括为长安失守后玄宗在蜀地难以直接统军，朔方将士与北方官僚需要一个能组织抗战的名义中心，并提示唐形成以灵武为军政中枢的复国政府，并为调动回纥等外援提供了可承认的君主。这类节点进入地方社会时，要经过官员、军队、书吏、商人、寺院、道路与家庭等不同中介；同一政策或战事在城市、边地和乡村不会具有相同的节奏。

核心来源为《旧唐书·肃宗本纪》；《新唐书·肃宗本纪》；《资治通鉴》唐纪（至德元载七月条）。它们可以较有把握地支持年序、制度公布、使节活动或特定地点的记载，却不能自动给出人口、财富、兵额、族属和内心动机的完整调查。账本的证据注记是“灵武即位的诏令、劝进程序和玄宗是否预先同意，主要由肃宗朝及后出的唐史保存”。因此，不能将官府的分类、战报的数字、宗教文本的愿望或后出的总括直接写成全体社会事实。

更合适的解释是，将这一事件看成资源、信息和权力重新连接的一环：命令是否抵达，取决于交通和地方协作；征发是否落实，取决于登记、市场和劳力；文化和宗教网络能否延续，也受战争、迁徙和赞助条件影响。现有材料只照亮其中一部分，叙事必须让区域差异与材料沉默继续可见。

条目\texttt{LT-756-003}所记的至德元载（756）“张巡在雍丘击退令狐潮部，守住汴宋间的据点”，发生在唐与燕的政治语境中。账本将其原因概括为燕军为控制运河与河南腹地进攻雍丘，张巡集合残余唐军和当地守备力量固守，并提示唐军在河南保留了一处牵制点，燕军南下补给线与对睢阳的作战安排受到拖延。这类节点进入地方社会时，要经过官员、军队、书吏、商人、寺院、道路与家庭等不同中介；同一政策或战事在城市、边地和乡村不会具有相同的节奏。

核心来源为《旧唐书·张巡传》；《新唐书·张巡传》；《资治通鉴》唐纪（至德元载条）。它们可以较有把握地支持年序、制度公布、使节活动或特定地点的记载，却不能自动给出人口、财富、兵额、族属和内心动机的完整调查。账本的证据注记是“张巡战绩、兵数和激烈言辞高度依赖忠臣传记，雍丘周边乡里与叛军材料几乎未存”。因此，不能将官府的分类、战报的数字、宗教文本的愿望或后出的总括直接写成全体社会事实。

更合适的解释是，将这一事件看成资源、信息和权力重新连接的一环：命令是否抵达，取决于交通和地方协作；征发是否落实，取决于登记、市场和劳力；文化和宗教网络能否延续，也受战争、迁徙和赞助条件影响。现有材料只照亮其中一部分，叙事必须让区域差异与材料沉默继续可见。

条目\texttt{LT-757-004}所记的至德二载（757）正月“安庆绪在洛阳杀安禄山，自立为帝”，发生在燕的政治语境中。账本将其原因概括为安禄山久病失明后猜忌近侍并威胁继承人，宫廷侍从与安庆绪在暴力威胁下结盟，并提示燕政权在战争高峰发生继承断裂，安庆绪对史思明等将领的控制随即减弱。这类节点进入地方社会时，要经过官员、军队、书吏、商人、寺院、道路与家庭等不同中介；同一政策或战事在城市、边地和乡村不会具有相同的节奏。

核心来源为《旧唐书·安禄山传》；《旧唐书·安庆绪传》；《资治通鉴》唐纪（至德二载正月条）。它们可以较有把握地支持年序、制度公布、使节活动或特定地点的记载，却不能自动给出人口、财富、兵额、族属和内心动机的完整调查。账本的证据注记是“弑父动机、严庄与李猪儿的责任分配主要见唐方敌国叙事，不能据此复原燕廷全部派系”。因此，不能将官府的分类、战报的数字、宗教文本的愿望或后出的总括直接写成全体社会事实。

更合适的解释是，将这一事件看成资源、信息和权力重新连接的一环：命令是否抵达，取决于交通和地方协作；征发是否落实，取决于登记、市场和劳力；文化和宗教网络能否延续，也受战争、迁徙和赞助条件影响。现有材料只照亮其中一部分，叙事必须让区域差异与材料沉默继续可见。

条目\texttt{LT-757-005}所记的至德二载（757）十月“尹子奇攻破睢阳，张巡、许远被俘后遇害”，发生在唐与燕的政治语境中。账本将其原因概括为睢阳守军兵少粮尽且外援迟至，尹子奇持续围城以夺取通往江淮的通路，并提示燕军虽占睢阳却付出长期围攻成本，江淮漕赋未立刻落入其手，忠烈记忆成为战后政治资源。这类节点进入地方社会时，要经过官员、军队、书吏、商人、寺院、道路与家庭等不同中介；同一政策或战事在城市、边地和乡村不会具有相同的节奏。

核心来源为《旧唐书·张巡传》；《旧唐书·许远传》；《资治通鉴》唐纪（至德二载十月条）。它们可以较有把握地支持年序、制度公布、使节活动或特定地点的记载，却不能自动给出人口、财富、兵额、族属和内心动机的完整调查。账本的证据注记是“守城饥荒和张巡、许远行为在忠烈叙事中被反复塑形，伤亡、食人等细节不可脱离文体核验”。因此，不能将官府的分类、战报的数字、宗教文本的愿望或后出的总括直接写成全体社会事实。

更合适的解释是，将这一事件看成资源、信息和权力重新连接的一环：命令是否抵达，取决于交通和地方协作；征发是否落实，取决于登记、市场和劳力；文化和宗教网络能否延续，也受战争、迁徙和赞助条件影响。现有材料只照亮其中一部分，叙事必须让区域差异与材料沉默继续可见。

条目\texttt{LT-757-006}所记的至德二载（757）九月“唐军与回纥骑兵在香积寺击败燕军，随后收复长安”，发生在唐与回纥、燕的政治语境中。账本将其原因概括为灵武政权需要机动力量突破关中战局，回纥则以婚盟、贸易和掠获回报作为结盟条件，并提示唐廷重获长安这一象征与行政中心，但对外援的酬谢和京城秩序恢复成为紧迫问题。这类节点进入地方社会时，要经过官员、军队、书吏、商人、寺院、道路与家庭等不同中介；同一政策或战事在城市、边地和乡村不会具有相同的节奏。

核心来源为《旧唐书·肃宗本纪》；《旧唐书·回纥传》；《资治通鉴》唐纪（至德二载九月条）。它们可以较有把握地支持年序、制度公布、使节活动或特定地点的记载，却不能自动给出人口、财富、兵额、族属和内心动机的完整调查。账本的证据注记是“香积寺战斗的兵数、回纥战利品与入城秩序多出自唐方记载，回纥行动逻辑须以其利益另作解释”。因此，不能将官府的分类、战报的数字、宗教文本的愿望或后出的总括直接写成全体社会事实。

更合适的解释是，将这一事件看成资源、信息和权力重新连接的一环：命令是否抵达，取决于交通和地方协作；征发是否落实，取决于登记、市场和劳力；文化和宗教网络能否延续，也受战争、迁徙和赞助条件影响。现有材料只照亮其中一部分，叙事必须让区域差异与材料沉默继续可见。

条目\texttt{LT-758-007}所记的乾元元年（758）“刘晏受命整顿盐铁转运，重建江淮财赋输送”，发生在唐的政治语境中。账本将其原因概括为安史战乱破坏河北、河南税源，肃宗朝需要从相对稳定的江淮组织粮饷和货币收入，并提示盐铁转运使成为连接中央与东南财赋的关键职位，财政重心进一步脱离传统租庸调。这类节点进入地方社会时，要经过官员、军队、书吏、商人、寺院、道路与家庭等不同中介；同一政策或战事在城市、边地和乡村不会具有相同的节奏。

核心来源为《旧唐书·刘晏传》；《新唐书·刘晏传》；《旧唐书·食货志下》。它们可以较有把握地支持年序、制度公布、使节活动或特定地点的记载，却不能自动给出人口、财富、兵额、族属和内心动机的完整调查。账本的证据注记是“刘晏任职的具体权限和改革先后在传记、食货志之间不尽一致，不能把后来的盐法成效前移到758年”。因此，不能将官府的分类、战报的数字、宗教文本的愿望或后出的总括直接写成全体社会事实。

更合适的解释是，将这一事件看成资源、信息和权力重新连接的一环：命令是否抵达，取决于交通和地方协作；征发是否落实，取决于登记、市场和劳力；文化和宗教网络能否延续，也受战争、迁徙和赞助条件影响。现有材料只照亮其中一部分，叙事必须让区域差异与材料沉默继续可见。

条目\texttt{LT-758-008}所记的乾元元年（758）九月“郭子仪等九节度使会攻相州邺城，唐军因内讧和史思明援军而溃”，发生在唐与燕的政治语境中。账本将其原因概括为唐廷急于消灭安庆绪而集中各镇军队，指挥权分散、补给竞争和将领猜忌削弱围城效率，并提示相州失利使唐军未能迅速结束战争，史思明得以借援军声望重建独立实力。这类节点进入地方社会时，要经过官员、军队、书吏、商人、寺院、道路与家庭等不同中介；同一政策或战事在城市、边地和乡村不会具有相同的节奏。

核心来源为《旧唐书·肃宗本纪》；《旧唐书·郭子仪传》；《资治通鉴》唐纪（乾元元年九月条）。它们可以较有把握地支持年序、制度公布、使节活动或特定地点的记载，却不能自动给出人口、财富、兵额、族属和内心动机的完整调查。账本的证据注记是“会攻部队数量与溃败责任被不同将领传记相互归咎，须区分战后自辩与编年叙事”。因此，不能将官府的分类、战报的数字、宗教文本的愿望或后出的总括直接写成全体社会事实。

更合适的解释是，将这一事件看成资源、信息和权力重新连接的一环：命令是否抵达，取决于交通和地方协作；征发是否落实，取决于登记、市场和劳力；文化和宗教网络能否延续，也受战争、迁徙和赞助条件影响。现有材料只照亮其中一部分，叙事必须让区域差异与材料沉默继续可见。

条目\texttt{LT-759-009}所记的乾元二年（759）三月“史思明杀安庆绪后据范阳、洛阳，自称大燕皇帝”，发生在唐与燕的政治语境中。账本将其原因概括为相州战后安庆绪兵败粮乏，史思明掌握范阳兵马并不愿受其猜忌和约束，并提示叛军由安氏父子继承转为史氏军事联盟，河北战场和唐军收复洛阳的计划再度受阻。这类节点进入地方社会时，要经过官员、军队、书吏、商人、寺院、道路与家庭等不同中介；同一政策或战事在城市、边地和乡村不会具有相同的节奏。

核心来源为《旧唐书·史思明传》；《新唐书·史思明传》；《资治通鉴》唐纪（乾元二年三月条）。它们可以较有把握地支持年序、制度公布、使节活动或特定地点的记载，却不能自动给出人口、财富、兵额、族属和内心动机的完整调查。账本的证据注记是“史思明与安庆绪决裂的经过主要由唐方传记串联，范阳军中支持者及地方财政状况缺少独立材料”。因此，不能将官府的分类、战报的数字、宗教文本的愿望或后出的总括直接写成全体社会事实。

更合适的解释是，将这一事件看成资源、信息和权力重新连接的一环：命令是否抵达，取决于交通和地方协作；征发是否落实，取决于登记、市场和劳力；文化和宗教网络能否延续，也受战争、迁徙和赞助条件影响。现有材料只照亮其中一部分，叙事必须让区域差异与材料沉默继续可见。

条目\texttt{LT-761-010}所记的上元二年（761）三月“史朝义发动政变杀史思明，继承燕军统帅权”，发生在燕的政治语境中。账本将其原因概括为史思明猜忌并威胁废立太子，史朝义为保全自身联合近侍和军中力量先发制人，并提示燕廷再次发生暴力继承，史朝义的号令难以覆盖各地将领，为763年瓦解埋下条件。这类节点进入地方社会时，要经过官员、军队、书吏、商人、寺院、道路与家庭等不同中介；同一政策或战事在城市、边地和乡村不会具有相同的节奏。

核心来源为《旧唐书·史思明传》；《旧唐书·史朝义传》；《资治通鉴》唐纪（上元二年三月条）。它们可以较有把握地支持年序、制度公布、使节活动或特定地点的记载，却不能自动给出人口、财富、兵额、族属和内心动机的完整调查。账本的证据注记是“弑父的直接过程与朝义同谋者只见敌对史书，燕军内部的日常命令链难以从现存材料复原”。因此，不能将官府的分类、战报的数字、宗教文本的愿望或后出的总括直接写成全体社会事实。

更合适的解释是，将这一事件看成资源、信息和权力重新连接的一环：命令是否抵达，取决于交通和地方协作；征发是否落实，取决于登记、市场和劳力；文化和宗教网络能否延续，也受战争、迁徙和赞助条件影响。现有材料只照亮其中一部分，叙事必须让区域差异与材料沉默继续可见。

条目\texttt{LT-762-011}所记的宝应元年（762）四月“唐肃宗卒，太子李豫即位，是为代宗”，发生在唐的政治语境中。账本将其原因概括为肃宗长期病重，李辅国控制禁军和诏令渠道，太子必须在皇位转换中获得军政支持，并提示新君接手的朝廷可继续协调复洛阳战役，但宦官干政和军功集团的权力问题更加突出。这类节点进入地方社会时，要经过官员、军队、书吏、商人、寺院、道路与家庭等不同中介；同一政策或战事在城市、边地和乡村不会具有相同的节奏。

核心来源为《旧唐书·肃宗本纪》；《旧唐书·代宗本纪》；《资治通鉴》唐纪（宝应元年四月条）。它们可以较有把握地支持年序、制度公布、使节活动或特定地点的记载，却不能自动给出人口、财富、兵额、族属和内心动机的完整调查。账本的证据注记是“肃宗末年宦官、后宫与太子关系的细节由宫廷政变叙事主导，死亡原因和即位程序有待慎读”。因此，不能将官府的分类、战报的数字、宗教文本的愿望或后出的总括直接写成全体社会事实。

更合适的解释是，将这一事件看成资源、信息和权力重新连接的一环：命令是否抵达，取决于交通和地方协作；征发是否落实，取决于登记、市场和劳力；文化和宗教网络能否延续，也受战争、迁徙和赞助条件影响。现有材料只照亮其中一部分，叙事必须让区域差异与材料沉默继续可见。

\subsection{边疆、道路与迁徙}

条目\texttt{LT-762-012}所记的宝应元年（762）十月“仆固怀恩引回纥军与唐军会师，进攻洛阳方向的史朝义军”，发生在唐、回纥与燕的政治语境中。账本将其原因概括为唐军欲乘燕内部继承危机收复两京，却缺少足够骑兵和机动作战能力，仆固怀恩推动再次借兵，并提示外援增强唐军攻势并加速燕军离散，同时使回纥在洛阳市场和唐廷财政中取得更大议价权。这类节点进入地方社会时，要经过官员、军队、书吏、商人、寺院、道路与家庭等不同中介；同一政策或战事在城市、边地和乡村不会具有相同的节奏。

核心来源为《旧唐书·仆固怀恩传》；《旧唐书·回纥传》；《资治通鉴》唐纪（宝应元年十月条）。它们可以较有把握地支持年序、制度公布、使节活动或特定地点的记载，却不能自动给出人口、财富、兵额、族属和内心动机的完整调查。账本的证据注记是“回纥出兵人数、掠夺和仆固怀恩意图在唐方史料中带有强烈政治指控，须同盟条件与战果分开考察”。因此，不能将官府的分类、战报的数字、宗教文本的愿望或后出的总括直接写成全体社会事实。

更合适的解释是，将这一事件看成资源、信息和权力重新连接的一环：命令是否抵达，取决于交通和地方协作；征发是否落实，取决于登记、市场和劳力；文化和宗教网络能否延续，也受战争、迁徙和赞助条件影响。现有材料只照亮其中一部分，叙事必须让区域差异与材料沉默继续可见。

条目\texttt{LT-763-013}所记的宝应二年（763）正月“史朝义兵败自缢，唐军重入洛阳，安史叛乱主战事结束”，发生在唐与燕的政治语境中。账本将其原因概括为史朝义失去部将支持且回纥、唐军联合压迫洛阳，各地燕将选择降唐或据地自保，并提示唐恢复两京的名义统一，却承认或容纳一批河北军镇，战后权力结构并未回到天宝旧制。这类节点进入地方社会时，要经过官员、军队、书吏、商人、寺院、道路与家庭等不同中介；同一政策或战事在城市、边地和乡村不会具有相同的节奏。

核心来源为《旧唐书·代宗本纪》；《旧唐书·史朝义传》；《资治通鉴》唐纪（宝应二年正月条）。它们可以较有把握地支持年序、制度公布、使节活动或特定地点的记载，却不能自动给出人口、财富、兵额、族属和内心动机的完整调查。账本的证据注记是“叛乱终结日期可由编年史确定，但河北诸镇的归附形式和实际税兵控制不能由洛阳收复推断”。因此，不能将官府的分类、战报的数字、宗教文本的愿望或后出的总括直接写成全体社会事实。

更合适的解释是，将这一事件看成资源、信息和权力重新连接的一环：命令是否抵达，取决于交通和地方协作；征发是否落实，取决于登记、市场和劳力；文化和宗教网络能否延续，也受战争、迁徙和赞助条件影响。现有材料只照亮其中一部分，叙事必须让区域差异与材料沉默继续可见。

条目\texttt{LT-763-014}所记的广德元年（763）十月“吐蕃军攻入长安，代宗出奔陕州，吐蕃旋即撤离”，发生在唐与吐蕃的政治语境中。账本将其原因概括为安史战后唐军主力仍在河南、河北，西北边防空虚，吐蕃利用唐廷与回纥关系紧张发起进攻，并提示长安再度失守显示复国并未恢复西部防线，唐朝不得不重新配置西北军费和外交资源。这类节点进入地方社会时，要经过官员、军队、书吏、商人、寺院、道路与家庭等不同中介；同一政策或战事在城市、边地和乡村不会具有相同的节奏。

核心来源为《旧唐书·代宗本纪》；《旧唐书·吐蕃传》；《资治通鉴》唐纪（广德元年十月条）。它们可以较有把握地支持年序、制度公布、使节活动或特定地点的记载，却不能自动给出人口、财富、兵额、族属和内心动机的完整调查。账本的证据注记是“吐蕃入长安的持续日数、所立傀儡和撤军原因在汉文与吐蕃材料中互有缺环，不能只作唐廷屈辱叙事”。因此，不能将官府的分类、战报的数字、宗教文本的愿望或后出的总括直接写成全体社会事实。

更合适的解释是，将这一事件看成资源、信息和权力重新连接的一环：命令是否抵达，取决于交通和地方协作；征发是否落实，取决于登记、市场和劳力；文化和宗教网络能否延续，也受战争、迁徙和赞助条件影响。现有材料只照亮其中一部分，叙事必须让区域差异与材料沉默继续可见。

条目\texttt{LT-764-015}所记的广德二年（764）“田承嗣据魏博并收纳安史旧将，形成相对独立的河北军镇”，发生在唐与魏博的政治语境中。账本将其原因概括为唐廷为尽快稳定河北而任用降将，田承嗣借旧部和地方仓储拒绝完全交还军政权，并提示河北藩镇的世袭化与自给性开始制度化，中央收复洛阳并不等于恢复对东部州县的直接控制。这类节点进入地方社会时，要经过官员、军队、书吏、商人、寺院、道路与家庭等不同中介；同一政策或战事在城市、边地和乡村不会具有相同的节奏。

核心来源为《旧唐书·田承嗣传》；《新唐书·田承嗣传》；《资治通鉴》唐纪（广德二年条）。它们可以较有把握地支持年序、制度公布、使节活动或特定地点的记载，却不能自动给出人口、财富、兵额、族属和内心动机的完整调查。账本的证据注记是“田承嗣控制州数和兵额主要来自中央史书与方镇表，魏博内部的地方服役、赋税结构仍待文献和考古补充”。因此，不能将官府的分类、战报的数字、宗教文本的愿望或后出的总括直接写成全体社会事实。

更合适的解释是，将这一事件看成资源、信息和权力重新连接的一环：命令是否抵达，取决于交通和地方协作；征发是否落实，取决于登记、市场和劳力；文化和宗教网络能否延续，也受战争、迁徙和赞助条件影响。现有材料只照亮其中一部分，叙事必须让区域差异与材料沉默继续可见。

条目\texttt{LT-765-016}所记的永泰元年（765）“唐与回纥重申绢马互市和婚姻、使节往来安排”，发生在唐与回纥的政治语境中。账本将其原因概括为唐廷需要稳定北方盟友并处理助战回纥的报酬，回纥则依赖唐绢与边市补充军政消费，并提示互市维系了脆弱同盟，也提高唐朝绢帛支出并使北边外交与财政政策更难分开。这类节点进入地方社会时，要经过官员、军队、书吏、商人、寺院、道路与家庭等不同中介；同一政策或战事在城市、边地和乡村不会具有相同的节奏。

核心来源为《旧唐书·代宗本纪》；《旧唐书·回纥传》；《唐会要》卷94〈北突厥〉。它们可以较有把握地支持年序、制度公布、使节活动或特定地点的记载，却不能自动给出人口、财富、兵额、族属和内心动机的完整调查。账本的证据注记是“会盟条款与马价在史书中多为唐廷视角，实际交易地点、商人群体和执行频率缺少连续记录”。因此，不能将官府的分类、战报的数字、宗教文本的愿望或后出的总括直接写成全体社会事实。

更合适的解释是，将这一事件看成资源、信息和权力重新连接的一环：命令是否抵达，取决于交通和地方协作；征发是否落实，取决于登记、市场和劳力；文化和宗教网络能否延续，也受战争、迁徙和赞助条件影响。现有材料只照亮其中一部分，叙事必须让区域差异与材料沉默继续可见。

条目\texttt{LT-766-017}所记的大历元年（766）“朱希彩杀李怀仙后自领卢龙军，朝廷予以承认”，发生在唐与卢龙的政治语境中。账本将其原因概括为李怀仙作为安史降将根基脆弱，朱希彩凭军中拥立夺权；唐廷无力在河北同时开战，并提示幽州的军人继承获得事实承认，卢龙成为唐后期长期自行决定主帅的北方镇区。这类节点进入地方社会时，要经过官员、军队、书吏、商人、寺院、道路与家庭等不同中介；同一政策或战事在城市、边地和乡村不会具有相同的节奏。

核心来源为《旧唐书·朱希彩传》；《新唐书·朱希彩传》；《资治通鉴》唐纪（大历元年条）。它们可以较有把握地支持年序、制度公布、使节活动或特定地点的记载，却不能自动给出人口、财富、兵额、族属和内心动机的完整调查。账本的证据注记是“朱希彩夺权的参与者和朝廷授任时间在纪传间略有出入，不能把诏授视为幽州军心一致服从”。因此，不能将官府的分类、战报的数字、宗教文本的愿望或后出的总括直接写成全体社会事实。

更合适的解释是，将这一事件看成资源、信息和权力重新连接的一环：命令是否抵达，取决于交通和地方协作；征发是否落实，取决于登记、市场和劳力；文化和宗教网络能否延续，也受战争、迁徙和赞助条件影响。现有材料只照亮其中一部分，叙事必须让区域差异与材料沉默继续可见。

条目\texttt{LT-768-018}所记的大历三年（768）“李正己兼并齐、棣等州，扩大平卢淄青镇的辖境”，发生在唐与淄青的政治语境中。账本将其原因概括为安史后平卢军南迁山东，唐廷对其节制薄弱，李正己趁周边州镇动荡扩展兵粮与辖区，并提示山东东部形成能控制沿海与运河支线的强镇，中央在该区的任官和征税空间进一步缩小。这类节点进入地方社会时，要经过官员、军队、书吏、商人、寺院、道路与家庭等不同中介；同一政策或战事在城市、边地和乡村不会具有相同的节奏。

核心来源为《旧唐书·李正己传》；《新唐书·李正己传》；《资治通鉴》唐纪（大历三年条）。它们可以较有把握地支持年序、制度公布、使节活动或特定地点的记载，却不能自动给出人口、财富、兵额、族属和内心动机的完整调查。账本的证据注记是“李正己所领州县和兼并时点在方镇表与本纪有不同表述，实际控制需与州县文献分别核查”。因此，不能将官府的分类、战报的数字、宗教文本的愿望或后出的总括直接写成全体社会事实。

更合适的解释是，将这一事件看成资源、信息和权力重新连接的一环：命令是否抵达，取决于交通和地方协作；征发是否落实，取决于登记、市场和劳力；文化和宗教网络能否延续，也受战争、迁徙和赞助条件影响。现有材料只照亮其中一部分，叙事必须让区域差异与材料沉默继续可见。

条目\texttt{LT-779-019}所记的大历十四年（779）五月“代宗卒，太子李适即位，是为德宗”，发生在唐的政治语境中。账本将其原因概括为代宗晚年依赖江淮财赋和对河北妥协，皇位继承发生在中央试图重整税兵关系之际，并提示德宗以削藩和财政整顿为目标的政策很快触发军镇、禁军与地方社会的连锁反应。这类节点进入地方社会时，要经过官员、军队、书吏、商人、寺院、道路与家庭等不同中介；同一政策或战事在城市、边地和乡村不会具有相同的节奏。

核心来源为《旧唐书·代宗本纪》；《旧唐书·德宗本纪上》；《资治通鉴》唐纪（大历十四年五月条）。它们可以较有把握地支持年序、制度公布、使节活动或特定地点的记载，却不能自动给出人口、财富、兵额、族属和内心动机的完整调查。账本的证据注记是“代宗遗命、宰相交接和即位礼细节以本纪为主，不能据其平稳叙事忽略财政和军镇压力”。因此，不能将官府的分类、战报的数字、宗教文本的愿望或后出的总括直接写成全体社会事实。

更合适的解释是，将这一事件看成资源、信息和权力重新连接的一环：命令是否抵达，取决于交通和地方协作；征发是否落实，取决于登记、市场和劳力；文化和宗教网络能否延续，也受战争、迁徙和赞助条件影响。现有材料只照亮其中一部分，叙事必须让区域差异与材料沉默继续可见。

条目\texttt{LT-780-020}所记的建中元年（780）“杨炎颁行两税法，以夏秋两次征税替代租庸调旧制”，发生在唐的政治语境中。账本将其原因概括为安史后户籍散失、均田租庸调难以征行，中央又需稳定军费，杨炎提出按资产与居住地征收，并提示税制从人丁与授田框架转向财产和资产核算的原则，但地方加派与藩镇留截仍限制中央收益。这类节点进入地方社会时，要经过官员、军队、书吏、商人、寺院、道路与家庭等不同中介；同一政策或战事在城市、边地和乡村不会具有相同的节奏。

核心来源为《旧唐书·杨炎传》；《新唐书·食货志五》；《通典》卷7〈食货七·历代盛衰户口〉。它们可以较有把握地支持年序、制度公布、使节活动或特定地点的记载，却不能自动给出人口、财富、兵额、族属和内心动机的完整调查。账本的证据注记是“两税法的诏令与后世志书概括不等于统一税率或即时执行，户等、土地和钱绢折纳须分地区检验”。因此，不能将官府的分类、战报的数字、宗教文本的愿望或后出的总括直接写成全体社会事实。

更合适的解释是，将这一事件看成资源、信息和权力重新连接的一环：命令是否抵达，取决于交通和地方协作；征发是否落实，取决于登记、市场和劳力；文化和宗教网络能否延续，也受战争、迁徙和赞助条件影响。现有材料只照亮其中一部分，叙事必须让区域差异与材料沉默继续可见。

条目\texttt{LT-781-021}所记的建中二年（781）“李惟岳等拒绝朝廷限制父死子继，河北多镇相继反叛”，发生在唐与成德、魏博等镇的政治语境中。账本将其原因概括为德宗试图以任命权收回藩镇继承，李惟岳、田悦等担忧失去家族与军中既得职位，并提示中央与河北军镇全面开战，德宗的削藩计划转而加重财政和对其他军队的依赖。这类节点进入地方社会时，要经过官员、军队、书吏、商人、寺院、道路与家庭等不同中介；同一政策或战事在城市、边地和乡村不会具有相同的节奏。

核心来源为《旧唐书·李惟岳传》；《旧唐书·田悦传》；《资治通鉴》唐纪（建中二年条）。它们可以较有把握地支持年序、制度公布、使节活动或特定地点的记载，却不能自动给出人口、财富、兵额、族属和内心动机的完整调查。账本的证据注记是““四镇之乱”的范围与各镇投入程度常被整合为单一叛乱，实际应区分成德、魏博、淄青与山南东道的行动”。因此，不能将官府的分类、战报的数字、宗教文本的愿望或后出的总括直接写成全体社会事实。

更合适的解释是，将这一事件看成资源、信息和权力重新连接的一环：命令是否抵达，取决于交通和地方协作；征发是否落实，取决于登记、市场和劳力；文化和宗教网络能否延续，也受战争、迁徙和赞助条件影响。现有材料只照亮其中一部分，叙事必须让区域差异与材料沉默继续可见。

条目\texttt{LT-783-022}所记的建中四年（783）十月“泾原军在长安哗变，朱泚据城称帝，德宗出奔奉天”，发生在唐与泾原军、朱泚的政治语境中。账本将其原因概括为河北战事抽空京畿资源，泾原军长期行军而待遇失衡，朱泚又拥有旧日京师军政网络，并提示德宗失去首都并被迫重新争取藩镇支持，削藩战争的政治目标由进攻转为自保。这类节点进入地方社会时，要经过官员、军队、书吏、商人、寺院、道路与家庭等不同中介；同一政策或战事在城市、边地和乡村不会具有相同的节奏。

核心来源为《旧唐书·德宗本纪上》；《旧唐书·朱泚传》；《资治通鉴》唐纪（建中四年十月条）。它们可以较有把握地支持年序、制度公布、使节活动或特定地点的记载，却不能自动给出人口、财富、兵额、族属和内心动机的完整调查。账本的证据注记是“泾原军哗变的赏赐、供给和市井冲突细节由不同叙事夸张，兵变不能被缩减为单一“赏赐不足”原因”。因此，不能将官府的分类、战报的数字、宗教文本的愿望或后出的总括直接写成全体社会事实。

更合适的解释是，将这一事件看成资源、信息和权力重新连接的一环：命令是否抵达，取决于交通和地方协作；征发是否落实，取决于登记、市场和劳力；文化和宗教网络能否延续，也受战争、迁徙和赞助条件影响。现有材料只照亮其中一部分，叙事必须让区域差异与材料沉默继续可见。

条目\texttt{LT-784-023}所记的兴元元年（784）“李希烈在蔡州称楚帝，公开脱离唐朝”，发生在唐与淮西的政治语境中。账本将其原因概括为德宗征调李希烈对抗河北诸镇却未能有效制衡其扩军，朝廷在奉天危机中缺乏制裁能力，并提示唐在东南、河南间的运输线受威胁，李希烈成为迫使朝廷重估地方军镇政策的又一中心。这类节点进入地方社会时，要经过官员、军队、书吏、商人、寺院、道路与家庭等不同中介；同一政策或战事在城市、边地和乡村不会具有相同的节奏。

核心来源为《旧唐书·李希烈传》；《新唐书·李希烈传》；《资治通鉴》唐纪（兴元元年条）。它们可以较有把握地支持年序、制度公布、使节活动或特定地点的记载，却不能自动给出人口、财富、兵额、族属和内心动机的完整调查。账本的证据注记是“李希烈称帝日期和国号较明，但其对淮西州县、士人和商路的实际控制不可由称号直接推定”。因此，不能将官府的分类、战报的数字、宗教文本的愿望或后出的总括直接写成全体社会事实。

更合适的解释是，将这一事件看成资源、信息和权力重新连接的一环：命令是否抵达，取决于交通和地方协作；征发是否落实，取决于登记、市场和劳力；文化和宗教网络能否延续，也受战争、迁徙和赞助条件影响。现有材料只照亮其中一部分，叙事必须让区域差异与材料沉默继续可见。

\subsection{文化、宗教与社群}

条目\texttt{LT-784-024}所记的兴元元年（784）正月“德宗发布奉天罪己诏，承诺赦免叛镇并调整政策”，发生在唐的政治语境中。账本将其原因概括为长安被朱泚占据、各镇叛乱扩大，德宗既缺兵力也缺政治信用，只能放弃先前强硬削藩路线，并提示部分藩镇转向勤王，朝廷获得复都的联盟条件，但以妥协换取效忠的模式被进一步固定。这类节点进入地方社会时，要经过官员、军队、书吏、商人、寺院、道路与家庭等不同中介；同一政策或战事在城市、边地和乡村不会具有相同的节奏。

核心来源为《旧唐书·德宗本纪上》；陆贽《奉天请罢琼林大盈二库状》；《资治通鉴》唐纪（兴元元年正月条）。它们可以较有把握地支持年序、制度公布、使节活动或特定地点的记载，却不能自动给出人口、财富、兵额、族属和内心动机的完整调查。账本的证据注记是“罪己诏的政治语言可确证，承诺是否转化为各镇具体待遇与民众感受须从后续诏令和地方材料检验”。因此，不能将官府的分类、战报的数字、宗教文本的愿望或后出的总括直接写成全体社会事实。

更合适的解释是，将这一事件看成资源、信息和权力重新连接的一环：命令是否抵达，取决于交通和地方协作；征发是否落实，取决于登记、市场和劳力；文化和宗教网络能否延续，也受战争、迁徙和赞助条件影响。现有材料只照亮其中一部分，叙事必须让区域差异与材料沉默继续可见。

条目\texttt{LT-785-025}所记的贞元元年（785）五月“德宗由奉天返回长安，恢复京师朝廷运作”，发生在唐的政治语境中。账本将其原因概括为李晟等收复长安、朱泚败亡，奉天朝廷得到藩镇和勤王军支持，具备回京条件，并提示中央重返都城却更依赖宦官、神策军与江淮财赋，京师恢复不意味着军镇问题已解决。这类节点进入地方社会时，要经过官员、军队、书吏、商人、寺院、道路与家庭等不同中介；同一政策或战事在城市、边地和乡村不会具有相同的节奏。

核心来源为《旧唐书·德宗本纪上》；《唐会要》卷30〈京城〉；《资治通鉴》唐纪（贞元元年五月条）。它们可以较有把握地支持年序、制度公布、使节活动或特定地点的记载，却不能自动给出人口、财富、兵额、族属和内心动机的完整调查。账本的证据注记是“还都礼仪可据本纪确定，但城中修复、居民回流和禁军重编的范围没有连续统计材料”。因此，不能将官府的分类、战报的数字、宗教文本的愿望或后出的总括直接写成全体社会事实。

更合适的解释是，将这一事件看成资源、信息和权力重新连接的一环：命令是否抵达，取决于交通和地方协作；征发是否落实，取决于登记、市场和劳力；文化和宗教网络能否延续，也受战争、迁徙和赞助条件影响。现有材料只照亮其中一部分，叙事必须让区域差异与材料沉默继续可见。

条目\texttt{LT-787-026}所记的贞元三年（787）“唐使在平凉会盟时遭吐蕃伏击，李晟等军队救出部分使者”，发生在唐与吐蕃的政治语境中。账本将其原因概括为唐在西北军费压力下寻求与吐蕃谈判，吐蕃则试图利用使团集结取得军事与外交优势，并提示盟约失败加深唐对陇右的防备，西北守军和对回纥、南诏的外交布局被迫调整。这类节点进入地方社会时，要经过官员、军队、书吏、商人、寺院、道路与家庭等不同中介；同一政策或战事在城市、边地和乡村不会具有相同的节奏。

核心来源为《旧唐书·吐蕃传》；《旧唐书·马燧传》；《资治通鉴》唐纪（贞元三年条）。它们可以较有把握地支持年序、制度公布、使节活动或特定地点的记载，却不能自动给出人口、财富、兵额、族属和内心动机的完整调查。账本的证据注记是“平凉会盟是否由双方同意的条约破裂、伏击责任与伤亡数在唐蕃叙事中差异很大，应保留外交文本缺失”。因此，不能将官府的分类、战报的数字、宗教文本的愿望或后出的总括直接写成全体社会事实。

更合适的解释是，将这一事件看成资源、信息和权力重新连接的一环：命令是否抵达，取决于交通和地方协作；征发是否落实，取决于登记、市场和劳力；文化和宗教网络能否延续，也受战争、迁徙和赞助条件影响。现有材料只照亮其中一部分，叙事必须让区域差异与材料沉默继续可见。

条目\texttt{LT-790-027}所记的贞元六年（790）“吐蕃控制北庭、安西余部的联络通道，唐在天山北路的军政体系终结”，发生在唐、吐蕃与北庭地区的政治语境中。账本将其原因概括为安史后唐廷撤回西北兵力，吐蕃持续东进并切断河西走廊，边镇难获中央补给，并提示唐朝对中亚的直接军事投射结束，丝路北线与西域地方社会转入多政权竞争的格局。这类节点进入地方社会时，要经过官员、军队、书吏、商人、寺院、道路与家庭等不同中介；同一政策或战事在城市、边地和乡村不会具有相同的节奏。

核心来源为《旧唐书·吐蕃传》；《新唐书·地理志七下》；《资治通鉴》唐纪（贞元六年条）。它们可以较有把握地支持年序、制度公布、使节活动或特定地点的记载，却不能自动给出人口、财富、兵额、族属和内心动机的完整调查。账本的证据注记是“北庭失守的确切时点、城镇归属与唐军残部去向见于零散汉文记载，须同敦煌和西域材料互校”。因此，不能将官府的分类、战报的数字、宗教文本的愿望或后出的总括直接写成全体社会事实。

更合适的解释是，将这一事件看成资源、信息和权力重新连接的一环：命令是否抵达，取决于交通和地方协作；征发是否落实，取决于登记、市场和劳力；文化和宗教网络能否延续，也受战争、迁徙和赞助条件影响。现有材料只照亮其中一部分，叙事必须让区域差异与材料沉默继续可见。

条目\texttt{LT-801-028}所记的贞元十七年（801）“韦皋在西川击退南诏、吐蕃联军，稳定成都西南防线”，发生在唐、西川、南诏与吐蕃的政治语境中。账本将其原因概括为吐蕃联结南诏以牵制唐西南，四川边军在韦皋长期经营下获得较稳定的城防和地方协作，并提示成都平原免于直接战火，西川节度使的军政地位和对南诏外交的影响随之上升。这类节点进入地方社会时，要经过官员、军队、书吏、商人、寺院、道路与家庭等不同中介；同一政策或战事在城市、边地和乡村不会具有相同的节奏。

核心来源为《旧唐书·韦皋传》；《新唐书·南蛮传下》；《资治通鉴》唐纪（贞元十七年条）。它们可以较有把握地支持年序、制度公布、使节活动或特定地点的记载，却不能自动给出人口、财富、兵额、族属和内心动机的完整调查。账本的证据注记是“韦皋传对战果和斩获多用报功数字，南诏、吐蕃联军的指挥与补给资料不足”。因此，不能将官府的分类、战报的数字、宗教文本的愿望或后出的总括直接写成全体社会事实。

更合适的解释是，将这一事件看成资源、信息和权力重新连接的一环：命令是否抵达，取决于交通和地方协作；征发是否落实，取决于登记、市场和劳力；文化和宗教网络能否延续，也受战争、迁徙和赞助条件影响。现有材料只照亮其中一部分，叙事必须让区域差异与材料沉默继续可见。

条目\texttt{LT-805-029}所记的永贞元年（805）“王叔文、王伾等在顺宗朝推动罢宫市、抑制宦官等永贞改革”，发生在唐的政治语境中。账本将其原因概括为顺宗即位后病重，王叔文集团借东宫与朝臣支持试图清理内廷经济和使职弊端，并提示改革迅速遭宦官和藩镇反制，晚唐政治中皇帝、宦官与文官集团的权力边界更加尖锐。这类节点进入地方社会时，要经过官员、军队、书吏、商人、寺院、道路与家庭等不同中介；同一政策或战事在城市、边地和乡村不会具有相同的节奏。

核心来源为《旧唐书·王叔文传》；《新唐书·王叔文传》；刘禹锡《子刘子自传》。它们可以较有把握地支持年序、制度公布、使节活动或特定地点的记载，却不能自动给出人口、财富、兵额、族属和内心动机的完整调查。账本的证据注记是“改革措施的颁行与实际废止次序需区分，王叔文集团形象深受宦官政变后传记与文学叙事影响”。因此，不能将官府的分类、战报的数字、宗教文本的愿望或后出的总括直接写成全体社会事实。

更合适的解释是，将这一事件看成资源、信息和权力重新连接的一环：命令是否抵达，取决于交通和地方协作；征发是否落实，取决于登记、市场和劳力；文化和宗教网络能否延续，也受战争、迁徙和赞助条件影响。现有材料只照亮其中一部分，叙事必须让区域差异与材料沉默继续可见。

条目\texttt{LT-805-030}所记的永贞元年（805）八月“顺宗让位太子李纯，宪宗即位，永贞改革集团被逐”，发生在唐的政治语境中。账本将其原因概括为顺宗病势加重，王叔文集团未能掌握神策军，俱文珍等宦官与太子周边官员促成权力移交，并提示宪宗获得较稳固的名义继承，王叔文集团退出中枢，削藩政策转入新的军事路径。这类节点进入地方社会时，要经过官员、军队、书吏、商人、寺院、道路与家庭等不同中介；同一政策或战事在城市、边地和乡村不会具有相同的节奏。

核心来源为《旧唐书·顺宗本纪》；《旧唐书·宪宗本纪上》；《资治通鉴》唐纪（永贞元年八月条）。它们可以较有把握地支持年序、制度公布、使节活动或特定地点的记载，却不能自动给出人口、财富、兵额、族属和内心动机的完整调查。账本的证据注记是“顺宗是否完全出于自愿让位及宦官介入程度，主要依赖宫廷史料和后世政治评价”。因此，不能将官府的分类、战报的数字、宗教文本的愿望或后出的总括直接写成全体社会事实。

更合适的解释是，将这一事件看成资源、信息和权力重新连接的一环：命令是否抵达，取决于交通和地方协作；征发是否落实，取决于登记、市场和劳力；文化和宗教网络能否延续，也受战争、迁徙和赞助条件影响。现有材料只照亮其中一部分，叙事必须让区域差异与材料沉默继续可见。

条目\texttt{LT-807-031}所记的元和二年（807）“刘辟据成都反叛，唐军入蜀后将其擒杀”，发生在唐与西川的政治语境中。账本将其原因概括为韦皋死后西川军镇权力交接不稳，刘辟拒绝朝廷另授节度使并依靠成都兵粮自立，并提示宪宗以较短战事收回西川，显示中央在非河北地区仍能调度军队和任命地方长官。这类节点进入地方社会时，要经过官员、军队、书吏、商人、寺院、道路与家庭等不同中介；同一政策或战事在城市、边地和乡村不会具有相同的节奏。

核心来源为《旧唐书·刘辟传》；《新唐书·刘辟传》；《资治通鉴》唐纪（元和二年条）。它们可以较有把握地支持年序、制度公布、使节活动或特定地点的记载，却不能自动给出人口、财富、兵额、族属和内心动机的完整调查。账本的证据注记是“刘辟起兵的动机、民众响应与唐军行军路线多取自中央军功叙事，四川地方社会资料十分稀少”。因此，不能将官府的分类、战报的数字、宗教文本的愿望或后出的总括直接写成全体社会事实。

更合适的解释是，将这一事件看成资源、信息和权力重新连接的一环：命令是否抵达，取决于交通和地方协作；征发是否落实，取决于登记、市场和劳力；文化和宗教网络能否延续，也受战争、迁徙和赞助条件影响。现有材料只照亮其中一部分，叙事必须让区域差异与材料沉默继续可见。

条目\texttt{LT-809-032}所记的元和四年（809）“王承宗拒还德、棣二州，宪宗用兵后暂予赦宥”，发生在唐与成德的政治语境中。账本将其原因概括为朝廷试图拆分成德镇并收回两州，王承宗依托镇兵抵抗，邻近河北军镇态度使战事难速决，并提示中央暂时承认成德现状，宪宗此后更重视选择可孤立、可切断补给的藩镇作为目标。这类节点进入地方社会时，要经过官员、军队、书吏、商人、寺院、道路与家庭等不同中介；同一政策或战事在城市、边地和乡村不会具有相同的节奏。

核心来源为《旧唐书·王承宗传》；《旧唐书·宪宗本纪上》；《资治通鉴》唐纪（元和四年条）。它们可以较有把握地支持年序、制度公布、使节活动或特定地点的记载，却不能自动给出人口、财富、兵额、族属和内心动机的完整调查。账本的证据注记是“战事中的州县归属、兵数和赦宥条件在本纪与传记中有差别，不能视作成德全境一致意愿”。因此，不能将官府的分类、战报的数字、宗教文本的愿望或后出的总括直接写成全体社会事实。

更合适的解释是，将这一事件看成资源、信息和权力重新连接的一环：命令是否抵达，取决于交通和地方协作；征发是否落实，取决于登记、市场和劳力；文化和宗教网络能否延续，也受战争、迁徙和赞助条件影响。现有材料只照亮其中一部分，叙事必须让区域差异与材料沉默继续可见。

条目\texttt{LT-813-033}所记的元和八年（813）“吴元济继领淮西，拒绝朝廷改授并袭扰邻州”，发生在唐与淮西的政治语境中。账本将其原因概括为吴少阳死后吴元济以军中拥立继承，宪宗拒绝承认其世袭并要求调整淮西辖区，并提示淮西问题由任命争议升级为跨年战争，成为检验宪宗削藩能力的核心战场。这类节点进入地方社会时，要经过官员、军队、书吏、商人、寺院、道路与家庭等不同中介；同一政策或战事在城市、边地和乡村不会具有相同的节奏。

核心来源为《旧唐书·吴元济传》；《新唐书·吴元济传》；《资治通鉴》唐纪（元和八年条）。它们可以较有把握地支持年序、制度公布、使节活动或特定地点的记载，却不能自动给出人口、财富、兵额、族属和内心动机的完整调查。账本的证据注记是“吴元济继镇和袭州的年序可核，但淮西内部支持结构主要见唐廷将相的奏议，缺少地方自述”。因此，不能将官府的分类、战报的数字、宗教文本的愿望或后出的总括直接写成全体社会事实。

更合适的解释是，将这一事件看成资源、信息和权力重新连接的一环：命令是否抵达，取决于交通和地方协作；征发是否落实，取决于登记、市场和劳力；文化和宗教网络能否延续，也受战争、迁徙和赞助条件影响。现有材料只照亮其中一部分，叙事必须让区域差异与材料沉默继续可见。

条目\texttt{LT-814-034}所记的元和九年（814）“宪宗正式发兵征讨吴元济，诸道军围攻蔡州周边”，发生在唐与淮西的政治语境中。账本将其原因概括为吴元济持续劫掠唐、邓等州，朝廷又认为淮西地处漕运要冲，不能继续姑息，并提示中央投入多镇军力围攻蔡州，战争迫使宪宗在宰相、监军和节度使之间重整指挥关系。这类节点进入地方社会时，要经过官员、军队、书吏、商人、寺院、道路与家庭等不同中介；同一政策或战事在城市、边地和乡村不会具有相同的节奏。

核心来源为《旧唐书·宪宗本纪上》；《旧唐书·裴度传》；《资治通鉴》唐纪（元和九年条）。它们可以较有把握地支持年序、制度公布、使节活动或特定地点的记载，却不能自动给出人口、财富、兵额、族属和内心动机的完整调查。账本的证据注记是“讨淮西的部队数、军费与战功常由各镇上报，实际参战规模和地方负担仍难量化”。因此，不能将官府的分类、战报的数字、宗教文本的愿望或后出的总括直接写成全体社会事实。

更合适的解释是，将这一事件看成资源、信息和权力重新连接的一环：命令是否抵达，取决于交通和地方协作；征发是否落实，取决于登记、市场和劳力；文化和宗教网络能否延续，也受战争、迁徙和赞助条件影响。现有材料只照亮其中一部分，叙事必须让区域差异与材料沉默继续可见。

条目\texttt{LT-815-035}所记的元和十年（815）六月“宰相武元衡在长安遇刺，御史中丞裴度亦遭袭”，发生在唐的政治语境中。账本将其原因概括为武元衡、裴度主张坚持征淮西，战争引发镇内外反对者与京城政治网络的尖锐冲突，并提示宪宗没有停战，裴度成为强硬派领袖，朝廷更以刺杀事件巩固继续用兵的政治正当性。这类节点进入地方社会时，要经过官员、军队、书吏、商人、寺院、道路与家庭等不同中介；同一政策或战事在城市、边地和乡村不会具有相同的节奏。

核心来源为《旧唐书·武元衡传》；《旧唐书·裴度传》；《资治通鉴》唐纪（元和十年六月条）。它们可以较有把握地支持年序、制度公布、使节活动或特定地点的记载，却不能自动给出人口、财富、兵额、族属和内心动机的完整调查。账本的证据注记是“刺客主使与淮西、成德的关联在审讯和后出史书中并非无疑，不能将政治指控当作已证事实”。因此，不能将官府的分类、战报的数字、宗教文本的愿望或后出的总括直接写成全体社会事实。

更合适的解释是，将这一事件看成资源、信息和权力重新连接的一环：命令是否抵达，取决于交通和地方协作；征发是否落实，取决于登记、市场和劳力；文化和宗教网络能否延续，也受战争、迁徙和赞助条件影响。现有材料只照亮其中一部分，叙事必须让区域差异与材料沉默继续可见。

\subsection{环境、财政与社会后果}

条目\texttt{LT-817-036}所记的元和十二年（817）十月“李愬夜袭蔡州，擒吴元济，淮西镇被平定”，发生在唐与淮西的政治语境中。账本将其原因概括为长期围攻消耗淮西粮食与人心，李愬通过招降李祐等人获得城防与行军情报，并提示宪宗取得最显著的削藩胜利，中央威信短暂回升，但后续接管仍依赖地方军人与财政资源。这类节点进入地方社会时，要经过官员、军队、书吏、商人、寺院、道路与家庭等不同中介；同一政策或战事在城市、边地和乡村不会具有相同的节奏。

核心来源为《旧唐书·李愬传》；《旧唐书·吴元济传》；《资治通鉴》唐纪（元和十二年十月条）。它们可以较有把握地支持年序、制度公布、使节活动或特定地点的记载，却不能自动给出人口、财富、兵额、族属和内心动机的完整调查。账本的证据注记是“夜袭路线、降雪与兵数在李愬传中具有英雄化叙事特征，蔡州内部倒戈的社会条件须另行追索”。因此，不能将官府的分类、战报的数字、宗教文本的愿望或后出的总括直接写成全体社会事实。

更合适的解释是，将这一事件看成资源、信息和权力重新连接的一环：命令是否抵达，取决于交通和地方协作；征发是否落实，取决于登记、市场和劳力；文化和宗教网络能否延续，也受战争、迁徙和赞助条件影响。现有材料只照亮其中一部分，叙事必须让区域差异与材料沉默继续可见。

条目\texttt{LT-819-037}所记的元和十四年（819）“宪宗迎法门寺佛骨入宫，韩愈上《论佛骨表》遭贬潮州”，发生在唐与佛教机构的政治语境中。账本将其原因概括为宪宗在削藩成功后寻求宗教性声望，法门寺舍利传统和内廷信奉使迎骨成为可行仪式，并提示儒佛争论被置入国家礼制议程，韩愈被贬成为后世讨论唐代宗教政策的标志性案例。这类节点进入地方社会时，要经过官员、军队、书吏、商人、寺院、道路与家庭等不同中介；同一政策或战事在城市、边地和乡村不会具有相同的节奏。

核心来源为《旧唐书·宪宗本纪下》；韩愈《论佛骨表》；《唐会要》卷47〈议释教上〉。它们可以较有把握地支持年序、制度公布、使节活动或特定地点的记载，却不能自动给出人口、财富、兵额、族属和内心动机的完整调查。账本的证据注记是“佛骨真伪与迎奉仪式的规模难由现存材料验证，韩愈奏表反映官僚论争而非全社会的宗教态度”。因此，不能将官府的分类、战报的数字、宗教文本的愿望或后出的总括直接写成全体社会事实。

更合适的解释是，将这一事件看成资源、信息和权力重新连接的一环：命令是否抵达，取决于交通和地方协作；征发是否落实，取决于登记、市场和劳力；文化和宗教网络能否延续，也受战争、迁徙和赞助条件影响。现有材料只照亮其中一部分，叙事必须让区域差异与材料沉默继续可见。

条目\texttt{LT-820-038}所记的元和十五年（820）正月“宪宗卒，穆宗李恒即位”，发生在唐的政治语境中。账本将其原因概括为宪宗晚年健康恶化且依赖宦官，太子已具继承资格但宫禁控制权集中于神策军系统，并提示穆宗和平继位，却较快转向对河北军镇的妥协，元和削藩的政治动能明显减退。这类节点进入地方社会时，要经过官员、军队、书吏、商人、寺院、道路与家庭等不同中介；同一政策或战事在城市、边地和乡村不会具有相同的节奏。

核心来源为《旧唐书·宪宗本纪下》；《旧唐书·穆宗本纪》；《资治通鉴》唐纪（元和十五年正月条）。它们可以较有把握地支持年序、制度公布、使节活动或特定地点的记载，却不能自动给出人口、财富、兵额、族属和内心动机的完整调查。账本的证据注记是“宪宗死因及陈弘志等人的角色在野史和正史中说法不一，不能把谋杀叙事写成定论”。因此，不能将官府的分类、战报的数字、宗教文本的愿望或后出的总括直接写成全体社会事实。

更合适的解释是，将这一事件看成资源、信息和权力重新连接的一环：命令是否抵达，取决于交通和地方协作；征发是否落实，取决于登记、市场和劳力；文化和宗教网络能否延续，也受战争、迁徙和赞助条件影响。现有材料只照亮其中一部分，叙事必须让区域差异与材料沉默继续可见。

条目\texttt{LT-821-039}所记的长庆元年（821）“礼部侍郎钱徽主持进士取士引发“牛李党争”早期科举争议”，发生在唐的政治语境中。账本将其原因概括为进士取士仍依赖荐举、文名和主考裁量，落第者与御史借榜单质疑钱徽的人事公正，并提示朝廷改易榜单并处罚相关官员，科举成为中晚唐党派语言和政治信誉竞争的重要场域。这类节点进入地方社会时，要经过官员、军队、书吏、商人、寺院、道路与家庭等不同中介；同一政策或战事在城市、边地和乡村不会具有相同的节奏。

核心来源为《旧唐书·钱徽传》；《新唐书·李宗闵传》；《唐会要》卷76〈贡举〉。它们可以较有把握地支持年序、制度公布、使节活动或特定地点的记载，却不能自动给出人口、财富、兵额、族属和内心动机的完整调查。账本的证据注记是“榜单舞弊指控和党争标签多由后出史论放大，考生名次、请托和派系归属不能逐一确证”。因此，不能将官府的分类、战报的数字、宗教文本的愿望或后出的总括直接写成全体社会事实。

更合适的解释是，将这一事件看成资源、信息和权力重新连接的一环：命令是否抵达，取决于交通和地方协作；征发是否落实，取决于登记、市场和劳力；文化和宗教网络能否延续，也受战争、迁徙和赞助条件影响。现有材料只照亮其中一部分，叙事必须让区域差异与材料沉默继续可见。

条目\texttt{LT-821-040}所记的长庆元年（821）“唐蕃订立长庆会盟，双方约定边界与使节往来，后立会盟碑”，发生在唐与吐蕃的政治语境中。账本将其原因概括为唐、吐蕃均受内政和边防消耗，双方需要稳定河陇边线并重开正式沟通渠道，并提示会盟暂时降低大规模冲突，留下罕见的碑刻文本，但未消除河西、青海与地方军镇的竞争。这类节点进入地方社会时，要经过官员、军队、书吏、商人、寺院、道路与家庭等不同中介；同一政策或战事在城市、边地和乡村不会具有相同的节奏。

核心来源为《旧唐书·穆宗本纪》；《旧唐书·吐蕃传》；《唐蕃会盟碑》（长庆三年立）。它们可以较有把握地支持年序、制度公布、使节活动或特定地点的记载，却不能自动给出人口、财富、兵额、族属和内心动机的完整调查。账本的证据注记是“会盟文现由长安、拉萨碑刻与汉藏史料保存，文本版本、边界地名和实际执行时间仍需逐项互校”。因此，不能将官府的分类、战报的数字、宗教文本的愿望或后出的总括直接写成全体社会事实。

更合适的解释是，将这一事件看成资源、信息和权力重新连接的一环：命令是否抵达，取决于交通和地方协作；征发是否落实，取决于登记、市场和劳力；文化和宗教网络能否延续，也受战争、迁徙和赞助条件影响。现有材料只照亮其中一部分，叙事必须让区域差异与材料沉默继续可见。

条目\texttt{LT-824-041}所记的长庆四年（824）正月“穆宗卒，敬宗李湛即位”，发生在唐的政治语境中。账本将其原因概括为穆宗长期病重，太子李湛已被册立，神策军宦官掌握宫禁并可影响仪式和诏令传达，并提示敬宗继位后宫廷消费与近侍政治受到批评，为文宗朝的内廷冲突留下背景。这类节点进入地方社会时，要经过官员、军队、书吏、商人、寺院、道路与家庭等不同中介；同一政策或战事在城市、边地和乡村不会具有相同的节奏。

核心来源为《旧唐书·穆宗本纪》；《旧唐书·敬宗本纪》；《资治通鉴》唐纪（长庆四年正月条）。它们可以较有把握地支持年序、制度公布、使节活动或特定地点的记载，却不能自动给出人口、财富、兵额、族属和内心动机的完整调查。账本的证据注记是“穆宗死因与敬宗早期朝政主要由宫廷史料书写，宦官、外戚和文官的实际影响需逐案区分”。因此，不能将官府的分类、战报的数字、宗教文本的愿望或后出的总括直接写成全体社会事实。

更合适的解释是，将这一事件看成资源、信息和权力重新连接的一环：命令是否抵达，取决于交通和地方协作；征发是否落实，取决于登记、市场和劳力；文化和宗教网络能否延续，也受战争、迁徙和赞助条件影响。现有材料只照亮其中一部分，叙事必须让区域差异与材料沉默继续可见。

条目\texttt{LT-828-042}所记的大和二年（828）“李同捷在沧州拒命起兵，宪宗后留下的镇将继承问题再起”，发生在唐与沧景的政治语境中。账本将其原因概括为李全略死后，李同捷仿效藩镇世袭要求承袭节度使，文宗朝廷拒绝追认其自立，并提示沧景战争迫使中央联络河北诸镇，显示在该地区用兵仍须依赖地方军镇的合作。这类节点进入地方社会时，要经过官员、军队、书吏、商人、寺院、道路与家庭等不同中介；同一政策或战事在城市、边地和乡村不会具有相同的节奏。

核心来源为《旧唐书·李同捷传》；《旧唐书·文宗本纪上》；《资治通鉴》唐纪（大和二年条）。它们可以较有把握地支持年序、制度公布、使节活动或特定地点的记载，却不能自动给出人口、财富、兵额、族属和内心动机的完整调查。账本的证据注记是“李同捷是否获得沧景军民广泛支持不能由朝廷檄文判断，战事叙述多以讨叛将领的军功为线索”。因此，不能将官府的分类、战报的数字、宗教文本的愿望或后出的总括直接写成全体社会事实。

更合适的解释是，将这一事件看成资源、信息和权力重新连接的一环：命令是否抵达，取决于交通和地方协作；征发是否落实，取决于登记、市场和劳力；文化和宗教网络能否延续，也受战争、迁徙和赞助条件影响。现有材料只照亮其中一部分，叙事必须让区域差异与材料沉默继续可见。

条目\texttt{LT-829-043}所记的大和三年（829）“魏博军擒李同捷，唐朝结束沧景之乱”，发生在唐与沧景的政治语境中。账本将其原因概括为李同捷孤立后粮道受阻，朝廷以授官、赏赐争取魏博何进滔等镇倒向讨伐一方，并提示横海镇被重新分置，唐廷获得一次有限胜利，却需继续以优待河北军镇换取地区合作。这类节点进入地方社会时，要经过官员、军队、书吏、商人、寺院、道路与家庭等不同中介；同一政策或战事在城市、边地和乡村不会具有相同的节奏。

核心来源为《旧唐书·李同捷传》；《新唐书·何进滔传》；《资治通鉴》唐纪（大和三年条）。它们可以较有把握地支持年序、制度公布、使节活动或特定地点的记载，却不能自动给出人口、财富、兵额、族属和内心动机的完整调查。账本的证据注记是“李同捷被擒的归属和魏博出兵动机带有朝廷嘉奖叙事，不能将其等同为河北诸镇重新受控”。因此，不能将官府的分类、战报的数字、宗教文本的愿望或后出的总括直接写成全体社会事实。

更合适的解释是，将这一事件看成资源、信息和权力重新连接的一环：命令是否抵达，取决于交通和地方协作；征发是否落实，取决于登记、市场和劳力；文化和宗教网络能否延续，也受战争、迁徙和赞助条件影响。现有材料只照亮其中一部分，叙事必须让区域差异与材料沉默继续可见。

条目\texttt{LT-835-044}所记的大和九年（835）十一月“李训、郑注谋诛宦官失败，甘露之变后宦官大杀朝臣”，发生在唐的政治语境中。账本将其原因概括为文宗、李训和郑注试图摆脱宦官控制，却未能将密谋者与京师军队置于同一指挥链，并提示宦官以神策军控制朝廷并清洗士大夫，晚唐皇帝与外朝官僚制衡内廷的空间大幅收缩。这类节点进入地方社会时，要经过官员、军队、书吏、商人、寺院、道路与家庭等不同中介；同一政策或战事在城市、边地和乡村不会具有相同的节奏。

核心来源为《旧唐书·文宗本纪下》；《旧唐书·李训传》；《资治通鉴》唐纪（大和九年十一月条）。它们可以较有把握地支持年序、制度公布、使节活动或特定地点的记载，却不能自动给出人口、财富、兵额、族属和内心动机的完整调查。账本的证据注记是“事件现场、暗号和死亡人数主要由宦官胜利后形成的材料与后出编年叙事保存，细节不可尽信”。因此，不能将官府的分类、战报的数字、宗教文本的愿望或后出的总括直接写成全体社会事实。

更合适的解释是，将这一事件看成资源、信息和权力重新连接的一环：命令是否抵达，取决于交通和地方协作；征发是否落实，取决于登记、市场和劳力；文化和宗教网络能否延续，也受战争、迁徙和赞助条件影响。现有材料只照亮其中一部分，叙事必须让区域差异与材料沉默继续可见。

条目\texttt{LT-840-045}所记的开成五年（840）“回鹘汗国为黠戛斯所破，大批回鹘部众向河西、河东等地迁徙”，发生在回鹘与唐北边的政治语境中。账本将其原因概括为内部继承冲突削弱回鹘汗国，黠戛斯南下进攻其牙帐，草原贸易网络随之失序，并提示唐失去长期北方盟友，边市、绢马贸易和河西的安全环境都出现新的流动人口与竞争者。这类节点进入地方社会时，要经过官员、军队、书吏、商人、寺院、道路与家庭等不同中介；同一政策或战事在城市、边地和乡村不会具有相同的节奏。

核心来源为《旧唐书·回鹘传》；《新唐书·回鹘传下》；《资治通鉴》唐纪（开成五年条）。它们可以较有把握地支持年序、制度公布、使节活动或特定地点的记载，却不能自动给出人口、财富、兵额、族属和内心动机的完整调查。账本的证据注记是“汗国覆灭的战事与迁徙规模由汉文、突厥语材料分别叙述，部众路线和身份不能视为单一整体”。因此，不能将官府的分类、战报的数字、宗教文本的愿望或后出的总括直接写成全体社会事实。

更合适的解释是，将这一事件看成资源、信息和权力重新连接的一环：命令是否抵达，取决于交通和地方协作；征发是否落实，取决于登记、市场和劳力；文化和宗教网络能否延续，也受战争、迁徙和赞助条件影响。现有材料只照亮其中一部分，叙事必须让区域差异与材料沉默继续可见。

条目\texttt{LT-841-046}所记的会昌元年（841）“唐廷在河东、振武等地处置来降回鹘部众，并限制其内徙”，发生在唐与回鹘部众的政治语境中。账本将其原因概括为回鹘汗国瓦解后部众抵近唐境，河东、振武守将既需利用骑兵又恐供给和治安失控，并提示安置政策使边防获得可用人力，同时加剧唐、回鹘余部与沙陀等群体间的资源竞争。这类节点进入地方社会时，要经过官员、军队、书吏、商人、寺院、道路与家庭等不同中介；同一政策或战事在城市、边地和乡村不会具有相同的节奏。

核心来源为《旧唐书·回鹘传》；《唐会要》卷94〈北突厥〉；《资治通鉴》唐纪（会昌元年条）。它们可以较有把握地支持年序、制度公布、使节活动或特定地点的记载，却不能自动给出人口、财富、兵额、族属和内心动机的完整调查。账本的证据注记是“唐廷对回鹘人数、安置地和威胁的判断带有边防焦虑，降附部众的自我选择难从汉文诏令中得见”。因此，不能将官府的分类、战报的数字、宗教文本的愿望或后出的总括直接写成全体社会事实。

更合适的解释是，将这一事件看成资源、信息和权力重新连接的一环：命令是否抵达，取决于交通和地方协作；征发是否落实，取决于登记、市场和劳力；文化和宗教网络能否延续，也受战争、迁徙和赞助条件影响。现有材料只照亮其中一部分，叙事必须让区域差异与材料沉默继续可见。

条目\texttt{LT-842-047}所记的会昌二年（842）“吐蕃赞普朗达玛遇弑，吐蕃高层继承冲突逐步扩大”，发生在吐蕃的政治语境中。账本将其原因概括为中央权力围绕宗教政策和继承安排积累矛盾，赞普遇弑使原有宫廷平衡崩解，并提示吐蕃难再维持对河西、青海的统一调度，为张议潮等河西地方势力争取空间创造条件。这类节点进入地方社会时，要经过官员、军队、书吏、商人、寺院、道路与家庭等不同中介；同一政策或战事在城市、边地和乡村不会具有相同的节奏。

核心来源为《新唐书·吐蕃传下》；《资治通鉴》唐纪（会昌二年条）；《贤者喜宴》所载朗达玛遇弑传述。它们可以较有把握地支持年序、制度公布、使节活动或特定地点的记载，却不能自动给出人口、财富、兵额、族属和内心动机的完整调查。账本的证据注记是“朗达玛死亡的时间、凶手与其后分裂节奏在汉文和后出藏文叙事中不一致，不能以单一年份概括吐蕃解体”。因此，不能将官府的分类、战报的数字、宗教文本的愿望或后出的总括直接写成全体社会事实。

更合适的解释是，将这一事件看成资源、信息和权力重新连接的一环：命令是否抵达，取决于交通和地方协作；征发是否落实，取决于登记、市场和劳力；文化和宗教网络能否延续，也受战争、迁徙和赞助条件影响。现有材料只照亮其中一部分，叙事必须让区域差异与材料沉默继续可见。
